\documentclass[11pt,a4paper]{article}
\usepackage[T2A]{fontenc}
\usepackage[utf8]{inputenc}
\usepackage[russian,english]{babel}
\usepackage{amsmath,amssymb,amsthm,mathtools}
\usepackage{setspace}
\usepackage{enumitem}
\usepackage[hidelinks]{hyperref}
\setstretch{1.3}
\frenchspacing
\hyphenpenalty=10000
\renewcommand{\leq}{\leqslant}
\renewcommand{\le}{\leqslant}
\renewcommand{\geq}{\geqslant}
\renewcommand{\ge}{\geqslant}
\newcommand{\head}[1]{\par\vspace{15pt}\noindent\textbf{\underline{#1}}\quad\ignorespaces}
\newcommand{\src}[1]{\unskip\nobreak\hfil\penalty50\hskip1em\hbox{}\nobreak\hfill(#1){\parfillskip=0pt\par}}
\begin{document}
\begin{center}
{\huge\bfseries Methods of proof -- 2\par}\vspace{0.8em}
{\Large Induction: strengthening, dimension, descent\par}\vspace{0.8em}
{\Large Homework\par}\vspace{0.8em}
{\foreignlanguage{russian}{13 сентября 2026 года}\par}
\end{center}\vspace{-3.3pt}

\noindent\textbf{\underline{Theoretical minimum.}}\quad Three reflexes. First: if the induction
step does not close, strengthen the statement~--- replace a constant by a function of $n$, add an
extra term, prove the inequality with a remainder. Second: choose the right parameter to induct
on~--- dimension, rank, degree, number of prime factors, size of a block~--- and if two parameters
are involved, order the pairs lexicographically. Third: over the integers a minimal counterexample
always exists, so descent is available; over the reals it is not, and the substitute is
compactness or a limiting argument. Zorn's lemma is the transfinite form of ``take a maximal
element'' and is needed exactly when the object cannot be built in countably many steps.

\medskip
\noindent\rule{\textwidth}{0.4pt}

\begin{enumerate}[itemsep=0.4em,topsep=7.4pt]
\item Prove that $n$ lines in general position divide the plane into $\frac{n^2+n+2}{2}$ parts,
and find and prove the analogous formula for hyperplanes in $\mathbb R^d$ by a double induction on
$n$ and $d$.

\item Prove by induction on the degree that a polynomial of degree $n$ with $n+1$ integer values
at consecutive integers has integer finite differences, and deduce that $\binom xk$ form a basis
of the integer-valued polynomials (compare with Polynomials -- 7).

\item Prove the Cauchy--Schwarz inequality by induction on the number of terms, and the AM--GM
inequality by Cauchy's forward-backward induction; explain why the second one needs a
non-standard induction scheme.

\item Prove that every finite group of order $p^k$ has a non-trivial centre, by induction on $k$
using the class equation, and deduce that a group of order $p^2$ is abelian.

\item Prove by descent that the equation $x^2+y^2+z^2=2xyz$ has no positive integer solutions, and
that $x^3+2y^3=4z^3$ has none either.

\item Prove by induction on the rank that every matrix over a field is equivalent to a matrix with
ones and zeros on the diagonal, and by induction on the size that a nilpotent matrix is similar
to a strictly upper triangular one.

\item Show that the statement ``every continuous function on $[0,1]$ is bounded'' becomes false
if $[0,1]$ is replaced by $[0,1]\cap\mathbb Q$, so that no argument using only finitely many
points at a time (such as an induction on a natural parameter) can prove it, and give the
compactness proof; then explain the general rule about when an induction on $\mathbb N$ is
available at all.

\item Prove by induction on the number of vertices that a planar graph has a vertex of degree at
most five, and deduce the five-colour theorem.

\item Prove the existence of a maximal ideal in every ring with unit and of a maximal linearly
independent subset in every vector space, each by Zorn's lemma; then show that for a countably
generated space the countable procedure suffices.

\item Prove Sperner's lemma by induction on the dimension, and deduce Brouwer's fixed point
theorem; compare with the one-dimensional case in Continuous functions -- 5.
\end{enumerate}
\end{document}
